% SOURCE_STATEMENT: "What values can delocalized eta invariants of Dirac operators on universal covers take in general; in particular, are they always algebraic numbers?"
% Preparation and final self-audit: CONTRIBUTING.md prompts C and D; source/audit artifacts retained in the local physics_statements_2026_10_11 batch.
% !TeX program = lualatex
\documentclass[11pt]{article}
\usepackage[margin=25mm]{geometry}
\usepackage{fontspec}
\setmainfont{Latin Modern Roman}
\usepackage{amsmath,amssymb,mathtools}
\usepackage{unicode-math}
\setmathfont{Latin Modern Math}
\usepackage{xurl,hyperref}
\hypersetup{colorlinks=true,urlcolor=blue}
\setcounter{secnumdepth}{0}
\setlength{\parindent}{0pt}
\setlength{\parskip}{5pt}
\setlength{\emergencystretch}{3em}
\newcommand{\sref}[2]{\hyperlink{source:#1:#2}{[S#2]}}
\newcommand{\cataloguescope}[1]{\paragraph{Recorded target.}#1}
\newcommand{\R}{\mathbb R}
\newcommand{\C}{\mathbb C}
\newcommand{\N}{\mathbb N}
\newcommand{\Z}{\mathbb Z}
\newcommand{\Q}{\mathbb Q}
\newcommand{\F}{\mathbb F}
\newcommand{\eps}{\varepsilon}
\begin{document}
% UnsolvedMath ID 30003894; source catalogue code OWR-16407-005
\setcounter{section}{30003893}
\section{Arithmetic Nature of Delocalized Eta Invariants}\label{30003894}
{\small UnsolvedMath ID 30003894 · Mathematical Physics}
\subsection{Definitions and mathematical statement}

\paragraph{Shared notation.}$\mathbb N=\{1,2,\ldots\}$; the standard number fields have their usual meanings. All additional conventions are stated locally.


\paragraph{Definitions and conventions.}
For an odd-dimensional Dirac operator $D$, an eta invariant measures spectral asymmetry through a regularized trace integral. On a universal cover with deck group $\Gamma$, a delocalized trace sums the integral kernel over a specified nontrivial conjugacy class of $\Gamma$. Replacing the usual trace by this trace gives a delocalized eta invariant when the summation and regularization converge. An algebraic number is a root of a nonzero polynomial with rational coefficients.

\paragraph{Target retained from the canonical source catalogue.}
\paragraph{Catalogue wording retained for traceability.}
What values can delocalized eta invariants of Dirac operators on universal covers take in general; in particular, are they always algebraic numbers?

\paragraph{Source scope and qualification.}
The group, conjugacy class, operator and convergence hypotheses are essential to define the invariant. The source does not ask for algebraicity of divergent or undefined expressions, and the general value question is broader than the highlighted algebraicity question.
\subsection{Short English statement}
Which numbers can delocalized spectral-asymmetry invariants take, and must they always be algebraic?
\subsection{Sources}
\begin{itemize}
\item[\hypertarget{source:30003894:1}{S1}] ULAM.AI and the UnsolvedMath Contributors, UnsolvedMath, record 30003894 (OWR-16407-005). Catalogue curation under CC BY 4.0; definitions and English exposition added locally.
\url{https://huggingface.co/datasets/ulamai/UnsolvedMath}.
\item[\hypertarget{source:30003894:2}{S2}] Non-commutative Geometry, Index Theory and Mathematical Physics (2018), Oberwolfach Reports, DOI 10.4171/owr/2018/32, pp. 1911--1981. Primary report accessed 2026-10-11; relevant PDF pages 12.
\url{https://doi.org/10.4171/owr/2018/32}.
\item[\hypertarget{source:30003894:3}{S3}] Xie--Yu, IMRN (catalogue-listed context; not independently consulted in this preparation).
\url{https://doi.org/10.1093/imrn/rnz170}.
\end{itemize}
\label{end:30003894}
\end{document}
